\section*{Chapitre \ref{ch:g2:subtraction} --- Soustraction}

\begin{solution}{exo:g2:subtraction:1}
$58 - 23 = 35$~; \quad $176 - 45 = 131$~; \quad $389 - 164 = 225$.
\end{solution}

\begin{solution}{exo:g2:subtraction:2}
$52 - 17 = 35$~; \quad $73 - 28 = 45$~; \quad $60 - 24 = 36$.
\end{solution}

\begin{solution}{exo:g2:subtraction:3}
$324 - 167 = 157$ (vérification $157 + 167 = 324$)~; \quad
$500 - 236 = 264$ (vérification $264 + 236 = 500$)~; \quad
$703 - 158 = 545$ (vérification $545 + 158 = 703$).
\end{solution}

\begin{solution}{exo:g2:subtraction:4}
$62 - 58 = 4$~; \quad $91 - 85 = 6$~; \quad $100 - 96 = 4$.
\end{solution}

\begin{solution}{exo:g2:subtraction:5}
$50 - 42 = 8$~; \quad $20 - 13 = 7$~; \quad $100 - 75 = 25$.
\end{solution}

\begin{solution}{exo:g2:subtraction:6}
L'an dernier~: $156 - 28 = 128$~cm.
\end{solution}

\begin{solution}{exo:g2:subtraction:7}
$85 - 39 = 46$~: vérification $46 + 39 = 85$ --- juste.
$214 - 68 = 156$~: vérification $156 + 68 = 224 \neq 214$ --- faux
(la bonne réponse est $146$).
\end{solution}

\begin{solution}{exo:g2:subtraction:8}
\guillemotleft{} s'envolent \guillemotright{}~: $47 - 19 = 28$.
\guillemotleft{} de plus cueillies \guillemotright{}~:
$47 + 19 = 66$. \guillemotleft{} combien de plus
\guillemotright{}~: $47 - 19 = 28$.
\end{solution}

\begin{solution}{exo:g2:subtraction:9}
De $45$ à $63$ en montant, c'est $18$, donc $63 - 18 = 45$~: le
trou est $18$.
\end{solution}

\begin{solution}{exo:g2:subtraction:10}
Remonter à partir de $58$~: avant d'enlever $18$, le nombre était
$58 + 18 = 76$~; avant d'ajouter $30$, c'était $76 - 30 = 46$. Le
nombre de départ était $46$ (vérification~:
$46 + 30 - 18 = 58$).
\end{solution}
